\chapter{Wyższe grupy homotopii i związek z homologią}
\label{ch:wyzsza-homotopia}
\index{grupa homotopii}

Grupa podstawowa z Rozdziału~\ref{ch:topologia-algebraiczna-podstawy}
rejestruje pętle. Wyższe grupy homotopii mierzą, czy można
ściągnąć do punktu odwzorowania całej sfery $S^n$. Homologia z
Rozdziału~\ref{ch:homologia-przestrzeni} bierze z kolei
zorientowane cykle i zapomina część informacji o sposobie ich
odwzorowania. Twierdzenie Hurewicza wskazuje dokładnie sytuację,
w której te dwa opisy mają ten sam pierwszy niezerowy wynik.
Twierdzenie Whiteheada mówi, kiedy wszystkie grupy homotopii
rozstrzygają typ homotopijny kompleksu CW.

W tym rozdziale pracujemy z przestrzeniami z wybranym punktem bazowym,
przede wszystkim z kompleksami CW. W argumentach komórkowych
wybieramy punkt bazowy będący wierzchołkiem. Gdy piszemy
$\pi_n(X)$ bez punktu bazowego, zakładamy spójność łukową i
wybieramy punkt bazowy; jego zmiana wzdłuż drogi daje izomorfizm, ale
sam wybór drogi może mieć znaczenie.

\section{Od pętli do odwzorowań sfer}
\label{sec:grupy-pi-n}

\begin{definicja}[Grupa $\pi_n$ i punkt bazowy]
\label{def:pi-n}
Niech $I=[0,1]$, $I^n=I\times\cdots\times I$, a
$\partial I^n$ niech oznacza sumę wszystkich jego ścian.
Dla przestrzeni z wybranym punktem $x_0$ kładziemy
\[
 \pi_n(X,x_0)
 =\big[(I^n,\partial I^n),(X,x_0)\big]_*
 =\big[(S^n,*),(X,x_0)\big]_*
 \qquad(n\geq1).
\]
Nawias oznacza klasy homotopii, w których
\emph{przez cały czas} brzeg kostki lub wyróżniony
punkt sfery jest wysyłany do $x_0$.
Utożsamienie pochodzi z $I^n/\partial I^n\cong S^n$.
Dla $n=1$ otrzymujemy $\pi_1(X,x_0)$ z
Definicji~\ref{def:droga-petla}. Symbol $\pi_0(X)$
oznacza jedynie zbiór składowych łukowych,
zwykle bez działania grupowego.
\end{definicja}

Obraz odwzorowania $S^2\to X$ można sobie wyobrażać jako
elastyczną powłokę rozpiętą w przestrzeni; jej
ściągalność do punktu wymaga ciągłej deformacji
\emph{całego odwzorowania}, nie tylko zniknięcia klasy
homologii reprezentowanej przez to odwzorowanie. Sfera $S^n$ jest brzegiem dysku $D^{n+1}$,
a także jest homeomorficzna z brzegiem kostki $I^{n+1}$.

\begin{twierdzenie}[Działanie i przemienność]
\label{thm:pi-n-abelowe}
Dla $n\geq1$ klasy w $\pi_n(X,x_0)$ tworzą grupę.
Dla $n\geq2$ grupa ta jest abelowa.
\end{twierdzenie}
\begin{proof}
Jeśli $f,g:(I^n,\partial I^n)\to(X,x_0)$,
przedstawmy ich sumę sklejeniem w pierwszej
współrzędnej:
\[
 (f+_1g)(s_1,\ldots,s_n)=
 \begin{cases}
 f(2s_1,s_2,\ldots,s_n),&0\leq s_1\leq\tfrac12,\\
 g(2s_1-1,s_2,\ldots,s_n),&\tfrac12\leq s_1\leq1.
 \end{cases}
\]
Na płaszczyźnie sklejenia oba wzory dają
$x_0$; tak samo jest na zewnętrznym brzegu.
Homotopie $f$ lub $g$ zachowujące brzeg
sklejają się tym samym wzorem, więc działanie
nie zależy od reprezentantów. Łączność
wynika z liniowej zmiany podziału pierwszej
współrzędnej na trzy odcinki, dokładnie jak
w dowodzie Twierdzenia~\ref{thm:grupa-podstawowa}.
Stałe odwzorowanie jest elementem neutralnym:
dodatkową połowę, na której odwzorowanie jest stałe,
można stopniowo zwęzić. Odbicie
$s_1\mapsto1-s_1$ reprezentuje odwrotność,
bo dwie przeciwne połówki ściągamy
równocześnie do punktu.

Gdy $n\geq2$, możemy niezależnie zdefiniować
drugie działanie $+_2$, sklejając w kierunku
$s_2$. Ma ono ten sam element neutralny
i tę samą konstrukcję homotopii. Podział
kwadratu $(s_1,s_2)$ na cztery ćwiartki
daje \emph{równość odwzorowań}, a nie tylko klas:
\[
 (a+_1 b)+_2(c+_1d)
   =(a+_2c)+_1(b+_2d).
\]
Teraz przechodzimy do klas homotopii i oznaczamy klasę stałą przez $0$.
Korzystając z neutralności $0$, otrzymujemy
\[
 a+_1b=(a+_2 0)+_1(0+_2 b)
      =(a+_1 0)+_2(0+_1 b)=a+_2b.
\]
Z kolei
\[
 a+_1b=(0+_2a)+_1(b+_2 0)
      =(0+_1b)+_2(a+_1 0)=b+_2a=b+_1a.
\]
Te ostatnie równości dotyczą klas; samo usuwanie stałych połówek
zmienia parametryzację odwzorowania.
To argument Eckmanna--Hiltona: dwa
zgodne sposoby sklejania wymuszają
przemienność. Użycie drugiej
współrzędnej jest niemożliwe dla
$n=1$, dlatego $\pi_1$ może być
nieabelowa.
\end{proof}

\begin{figure}[!htbp]
\centering
\begin{tikzpicture}[>=Stealth,line cap=round,line join=round]
 \fill[manifoldblue!8] (0,0) rectangle (4,4);
 \draw[manifoldblue!75!black,line width=1.4pt] (0,0) rectangle (4,4);
 \draw[manifoldblue!55] (2,0)--(2,4) (0,2)--(4,2);
 \fill[manifoldred!20] (.28,2.28) rectangle (1.72,3.72);
 \fill[manifoldgreen!22] (2.28,.28) rectangle (3.72,1.72);
 \node[font=\large,manifoldred!80!black] at (1,3) {$f$};
 \node[font=\large,manifoldgreen!70!black] at (3,1) {$g$};
 \node at (3,3) {$*$}; \node at (1,1) {$*$};
 \draw[->,manifoldgold!85!black,thick] (0,-.35)--(4,-.35)
  node[midway,below=4pt] {sklejanie w kierunku $s_1$};
 \draw[->,manifoldgold!85!black,thick] (-.35,0)--(-.35,4)
  ;
 \node[rotate=90,manifoldgold!85!black] at (-.8,2)
  {sklejanie w kierunku $s_2$};
 \node[font=\small,align=left,anchor=west,text width=4.2cm]
  at (4.6,2)
  {Puste pola są odwzorowaniem stałym. Oba kierunki
   sklejania opisują to samo odwzorowanie czterech
   ćwiartek; zamiana kolejności daje
   $[f]+[g]=[g]+[f]$.};
\end{tikzpicture}
\caption{Dwa niezależne kierunki sklejania w $\pi_2$.
W $\pi_n$ dla $n\geq2$ pozostałe współrzędne
traktujemy jako parametry.}
\label{fig:pi2-sklejanie}
\end{figure}

\begin{stwierdzenie}[Odwzorowania, iloczyny i homotopie]
\label{prop:pi-n-funktor}
Ciągłe odwzorowanie oparte $F:(X,x_0)\to(Y,y_0)$ wyznacza
homomorfizm $F_*:\pi_n(X,x_0)\to\pi_n(Y,y_0)$
przez $[f]\mapsto[F\circ f]$.
Odwzorowania homotopijne \emph{przez odwzorowania oparte}
indukują ten sam homomorfizm. Dla
$n\geq1$ i przestrzeni opartych $X,Y$:
\[
 \pi_n(X\times Y,(x_0,y_0))
 \cong \pi_n(X,x_0)\times\pi_n(Y,y_0).
\]
\end{stwierdzenie}
\begin{proof}
Złożenie $F\circ f$ nadal wysyła
$\partial I^n$ do $y_0$, a homotopia
$f_t$ daje homotopię $F\circ f_t$.
Ponieważ $F$ działa po sklejaniu
obu połówek dokładnie tak samo
jak przed nim, $F_*([f]+[g])
=F_*[f]+F_*[g]$. Homotopia
oparta $F_t$ daje homotopię
$F_t\circ f$ z nieruchomym
obrazem brzegu. Wreszcie odwzorowanie
$I^n\to X\times Y$ to dokładnie
para odwzorowań do $X$ i $Y$; to
samo dotyczy homotopii.
Projekcje i składanie pary
stanowią wzajemne odwrotności
podanego izomorfizmu.
\end{proof}

\begin{uwaga}[Zmiana punktu bazowego]
Droga $\gamma$ od $x_0$ do $x_1$
pozwala umieścić odwzorowanie oparte w $x_1$
w mniejszej kostce i w otaczającym
ją pasie doprowadzić wartości
wzdłuż $\gamma$ do $x_0$.
Otrzymujemy izomorfizm grup
$\pi_n(X,x_1)\to\pi_n(X,x_0)$
z odwrotnością od drogi
przebytej w przeciwną stronę.
Dwie różne drogi mogą dać
różne izomorfizmy; dla zamkniętej
drogi dostajemy \emph{działanie}
$\pi_1(X,x_0)$ na $\pi_n(X,x_0)$.
Konstrukcja zależy tylko od klasy homotopii drogi względem końców:
homotopię drogi stosujemy na tym samym pasie. Złożenie dwóch pasów
odpowiada złożeniu dróg, a droga i jej odwrotność dają pas ściągalny.
To uzasadnia zarówno izomorfizm, jak i działanie grupy.
W przypadku przestrzeni jednospójnej wybór drogi
nie powoduje tej niejednoznaczności.
Jeśli $F_t:X\to Y$ jest homotopią bez warunku oparcia i
$\gamma(t)=F_t(x_0)$, to po zmianie punktu bazowego wzdłuż $\gamma$
homomorfizm $(F_1)_*$ staje się $(F_0)_*$: pas odwzorowania
$F_t\circ f$ jest właśnie pasem tej zmiany bazy.
W szczególności zwykła równoważność homotopijna indukuje
izomorfizmy wszystkich grup homotopii, z odpowiednio dobranymi bazami.
\end{uwaga}

\section{Grupy względne i długi ciąg homotopii}
\label{sec:pi-wzgledne}

Tak jak $H_n(X,A)$ mierzy cykle,
których brzeg może trafić do $A$,
tak względna homotopia mierzy
dyski w $X$ z brzegiem w $A$.
Brzeg dysku może się przy tym
przesuwać w $A$ podczas homotopii.

\begin{definicja}[Względne grupy homotopii]
\index{wzgledne grupy homotopii@Względne grupy homotopii}
\label{def:pi-wzgledne}
Niech $x_0\in A\subseteq X$.
W kostce $I^n$ wyróżnijmy
dolną ścianę $B=I^{n-1}\times\{0\}$,
a przez $J$ oznaczmy sumę
pozostałych ścian kostki $I^n$.
Klasa względna jest klasą odwzorowań
\[
 F:(I^n,\partial I^n,J)
       \longrightarrow(X,A,x_0);
\]
zatem $F(\partial I^n)\subseteq A$
i $F(J)=x_0$. Homotopie spełniają
te same warunki. Ich zbiór oznaczamy
$\pi_n(X,A,x_0)$.
W równoważnym modelu są to
oparte odwzorowania trójek
$(D^n,S^{n-1},*)\to(X,A,x_0)$.
Dla $n\geq2$ składanie w pierwszej
współrzędnej daje grupę; dla
$n\geq3$ jest ona abelowa.
Dla $n=1$ pozostaje \emph{zbiór
z wyróżnionym elementem},
ponieważ nie ma wolnej współrzędnej
do składania.
\end{definicja}

W przeciwieństwie do $\pi_n(X)$,
względna $\pi_2(X,A)$ może być
nieprzemienna. Drugi kierunek
sklejania zajmuje wówczas rolę
wyróżnionej ściany, a argument
z Rysunku~\ref{fig:pi2-sklejanie}
nie ma zastosowania.

\begin{figure}[!htbp]
\centering
\begin{tikzpicture}[>=Stealth,line cap=round,line join=round]
 \filldraw[rounded corners=8pt,fill=manifoldblue!7,draw=manifoldblue!60!black]
  (0,0) rectangle (8,3.4);
 \fill[manifoldblue!18] (2.5,1.65) ellipse (1.65 and 1.05);
 \draw[manifoldgreen!28,line width=12pt] (2.5,1.65) ellipse (1.65 and 1.05);
 \draw[manifoldred,line width=1.4pt] (2.5,1.65) ellipse (1.65 and 1.05);
 \fill[manifoldgold!80!black] (.85,1.65) circle(2.6pt);
 \node[left,font=\small] at (.78,1.65) {$x_0$};
 \node[manifoldblue!75!black] at (2.5,1.65) {$F(\operatorname{int}D^2)$};
 \node[anchor=west,manifoldgreen!60!black] at (5.1,2.5) {$A$ (zielony pas)};
 \draw[->,manifoldgreen!60!black] (5,2.5)--(3.85,2.32);
 \node[anchor=west,manifoldred!80!black] at (5.1,1.05) {$F(S^1)\subset A$};
 \draw[->,manifoldred!80!black] (5,1.05)--(4.05,1.3);
 \node[manifoldblue!75!black] at (7.65,3.05) {$X$};
\end{tikzpicture}
\caption{Przykład odwzorowania względnego $(D^2,S^1,*)\to(X,A,x_0)$.
Czerwony obraz brzegu leży w zielonym pasie $A$, podczas gdy
obraz wnętrza dysku może wyjść poza $A$. W ogólnej definicji
nie wymaga się różnowartościowości $F$.}
\label{fig:pi-wzgledny-dysk}
\end{figure}

\begin{twierdzenie}[Długi ciąg homotopii pary]
\label{thm:les-pi}
Dla $x_0\in A\subseteq X$ otrzymujemy
naturalny ciąg
\begin{equation}\label{eq:les-pi}
 \cdots\to\pi_n(A,x_0)\xrightarrow{i_*}
 \pi_n(X,x_0)\xrightarrow{j_*}
 \pi_n(X,A,x_0)\xrightarrow{\partial}
 \pi_{n-1}(A,x_0)\xrightarrow{i_*}
 \pi_{n-1}(X,x_0)\to\cdots .
\end{equation}
Jest on dokładny: obraz każdej
strzałki to jądro następnej.
Odwzorowanie $\partial$ ogranicza
reprezentujący dysk do jego
wyróżnionej ściany $B$.
Na końcu ciągu występują
zbiory $\pi_1(X,A,x_0)$ i
$\pi_0(A),\pi_0(X)$; dokładność
oznacza tam przeciwobraz
wyróżnionego elementu.
\end{twierdzenie}
\begin{proof}
W modelu kostkowym $i_*$ jest
włączeniem $A\to X$, a $j_*$
traktuje odwzorowanie z całym brzegiem
w $x_0$ jako odwzorowanie względne.
Wzór $\partial[F]=[F|_B]$
jest poprawny, ponieważ
$\partial B\subset J$ i
$F(J)=x_0$. Homotopia
względna ogranicza się
do homotopii opartych
odwzorowań na $B$, więc
$\partial$ nie zależy
od reprezentanta.

Sprawdźmy trzy kolejne miejsca.
\emph{W $\pi_n(X)$.} Jeśli $f$ ma cały obraz w $A$, homotopia
$f_t(s',r)=f(s',(1-t)r+t)$ jest względną homotopią do stałej;
tutaj $s'\in I^{n-1}$, a $r$ jest ostatnią współrzędną.
Odwrotnie, niech $H(s',r,t)$ będzie względną homotopią
od $f$ do stałej. Ściana $g(s',t)=H(s',0,t)$ leży w $A$
i ma cały brzeg w $x_0$, więc reprezentuje element $\pi_n(A)$.
W kwadracie $(r,t)$ porównujemy drogi od $(0,0)$ do $(1,1)$:
dolny, a potem prawy bok oraz lewy, a potem górny bok.
Są homotopijne względem końców przez liniową interpolację w kwadracie.
Po złożeniu z $H$ dają odpowiednio $f$ i $g$ ze stałym dodatkowym
odcinkiem; na brzegu współrzędnych $s'$ wszystko pozostaje w $x_0$.
Stąd $[f]=i_*[g]$.

\emph{W $\pi_n(X,A)$.} Dysk
pochodzący z $\pi_n(X)$
ma stały brzeg, więc jego
$\partial$ jest zerem.
Jeśli natomiast brzeg
$F|_B:S^{n-1}\to A$ jest
zerowy w $\pi_{n-1}(A)$,
wybieramy jego wypełnienie
przez $D^n\to A$. Sklejenie
tego wypełnienia z $F$
wzdłuż wspólnego brzegu
daje oparte odwzorowanie
$S^n=D^n\cup_{S^{n-1}}D^n\to X$.
Po zapomnieniu drugiego
dysku jego obraz względny
jest klasą $[F]$.

\emph{W $\pi_{n-1}(A)$.} Brzeg
dysku względnego wypełnia
się w $X$, więc ginie
po włączeniu do $X$.
Jeśli oparte odwzorowanie
$S^{n-1}\to A$ jest
zerowe w $\pi_{n-1}(X)$,
jego homotopia do odwzorowania
stałego wyznacza wypełnienie
$D^n\to X$ i tym samym
klasę względną, której
brzeg jest daną klasą.
Te trzy rachunki powtarzają się w sąsiednich stopniach.
Na końcu ciągu element $\pi_1(X,A,x_0)$ jest drogą
od pewnego $a\in A$ do $x_0$, a $\partial$ przypisuje mu
składową $a$ w $A$. Trafia ona do wyróżnionej składowej $X$.
Odwrotnie, punkt tej składowej można połączyć drogą z $x_0$.
Jeśli $a$ leży w składowej $x_0$ już w $A$, dopinamy drogę
z $x_0$ do $a$ w $A$ i otrzymujemy pętlę w $X$ reprezentującą
daną klasę względną. To dowodzi dokładności także dla zbiorów wskazanych.
Wszystkie konstrukcje zachowują się przy składaniu z mapą par,
co daje naturalność ciągu.
\end{proof}

\begin{definicja}[Spójność homotopijna]
\index{spojnosc homotopijna@Spójność homotopijna}
\label{def:spojnosc-homotopijna}
Przestrzeń $X$ jest
\emph{$r$-spójna} ($r\geq0$),
jeśli jest łukowo spójna
i $\pi_i(X)=0$ dla $1\leq i\leq r$.
W szczególności $1$-spójna
oznacza jednospójną.
Para $(X,A)$ jest
\emph{$r$-spójna}, jeśli
każda składowa łukowa $X$
spotyka $A$ oraz
$\pi_i(X,A)=0$ dla
$1\leq i\leq r$ przy każdym wyborze punktu bazowego w $A$
(w stopniu $1$ jest to
trywialność zbioru
z wyróżnioną klasą).
\end{definicja}

Dokładność~\eqref{eq:les-pi}
pokazuje, że dla $r$-spójnej
pary włączenie $A\hookrightarrow X$
jest izomorfizmem na
$\pi_i$ dla $i<r$ oraz
surjekcją na $\pi_r$.
Wymiar komórek CW pozwoli
za chwilę uzyskać takie
zanikanie bez liczenia odwzorowań
pojedynczo.

\section{Homotopie na komórkach CW}
\label{sec:pi-n-cw}

W rozdziale o homologii komórki dawały
generatory grup łańcuchów. W homotopii
ich wymiary pozwalają \emph{usuwać
przeszkody do deformacji}. Żeby
deformować jedną komórkę bez psucia
tego, co zrobiliśmy na jej brzegu,
potrzebujemy przedłużania homotopii.

\begin{lemat}[Zwarte obrazy i homotopie w CW]
\label{lem:zwartosc-cw-pi}
Każdy zwarty podzbiór kompleksu CW leży w skończonym podkompleksie.
Ponadto odwzorowanie $H:X\times I\to Y$ jest ciągłe, jeśli jego złożenia
z każdym $\chi_\alpha\times\id_I:D^{d_\alpha}\times I\to X\times I$
są ciągłe, gdzie $\chi_\alpha$ są mapami charakterystycznymi komórek.
\end{lemat}
\begin{proof}
Gdyby zwarty zbiór $C$ spotykał nieskończenie wiele otwartych komórek,
wybralibyśmy z każdej z nich jeden punkt $c_\alpha\in C$.
Każda domknięta komórka spotyka tylko skończenie wiele spośród tych punktów.
Każdy podzbiór $E$ zbioru wybranych punktów ma więc domknięty
przeciwobraz przez każdą mapę charakterystyczną: jest to skończona
suma przeciwobrazów punktów w przestrzeni Hausdorffa.
Topologia słaba mówi, że $E$ jest domknięty w $X$.
Zatem zbiór wszystkich $c_\alpha$ jest domkniętym, dyskretnym
podzbiorem zwartego $C$, co jest niemożliwe, jeśli jest nieskończony.
Do skończenie wielu spotkanych komórek dołączamy ich ściany;
warunek skończoności domknięć i malejące wymiary dają skończony podkompleks.

Topologia słaba oznacza, że surjekcja
$q:\bigsqcup_\alpha D^{d_\alpha}\to X$ jest ilorazowa.
Iloczyn dowolnej mapy ilorazowej $q:E\to X$ z $\id_I$ też jest ilorazowy.
Aby to sprawdzić, niech przeciwobraz $W\subseteq X\times I$ będzie otwarty.
Dla $(x,t)\in W$ wybieramy zwarte otoczenie $J$ punktu $t$ w $I$
takie, że $\{x\}\times J\subset W$. Zbiór
$V=\{y:\{y\}\times J\subset W\}$ ma otwarty przeciwobraz przez $q$:
dla każdego $e$ nad $V$ skończone pokrycie zwartego $J$
prostokątami zawartymi w przeciwobrazie $W$ daje otoczenie $e$
o tej samej własności. Stąd $V$ jest otwarty i
$V\times\operatorname{int}_I J\subseteq W$ jest otoczeniem $(x,t)$.
To dowodzi tezy o iloczynie, a więc także kryterium ciągłości $H$.
\end{proof}

Z lematu wynika również rozszerzenie homologii komórkowej
z Rozdziału~\ref{ch:homologia-przestrzeni} na dowolne CW i pary CW.
Dla pary skończonych CW definiujemy najpierw
$C_\bullet^{\mathrm{CW}}(X,A)=C_\bullet^{\mathrm{CW}}(X)/C_\bullet^{\mathrm{CW}}(A)$:
jest to kompleks wolny na komórkach spoza $A$, z brzegiem modulo komórki $A$.
Krótki ciąg tych kompleksów daje długi ciąg homologii.
Naturalne porównanie ze zwykłym ciągiem pary i lemat pięciu
przenoszą izomorfizmy dla $A,X$ na grupy względne.
Każdy singularny łańcuch, jego brzeg oraz świadectwo homologiczności
są skończonymi sumami map zwartych sympleksów, więc mieszczą się
w skończonym podkompleksie. Cykl pochodzi z takiego podkompleksu,
a relacja między cyklami jest już prawdziwa w pewnym większym,
również skończonym. Tak samo jest dla skończonych sum komórek
i dla par: przecięcie skończonego podkompleksu z podkompleksem $A$
jest skończonym podkompleksem. Izomorfizmy dla modeli skończonych
są zgodne z włączeniami, bo pochodzą z naturalnych ciągów par.
Przechodzą więc na izomorfizmy dla całego CW.
Grupy łańcuchów komórkowych są przy tym \emph{sumami prostymi},
a nie iloczynami: każda kombinacja nadal ma skończony nośnik.

\begin{lemat}[Przedłużanie homotopii dla par CW]
\label{lem:hep-cw}
Niech $A\subseteq X$ będzie podkompleksem
CW. Jeśli $f:X\to Y$ i homotopia
$h:A\times I\to Y$ spełniają
$h(a,0)=f(a)$, to istnieje
$H:X\times I\to Y$ takie, że
$H(x,0)=f(x)$ oraz $H|_{A\times I}=h$.
W szczególności homotopię
zadaną na brzegu dołączonego
dysku można przedłużyć do
całego dysku.
\end{lemat}
\begin{proof}
Najpierw rozważmy jedną komórkę
$D^k$, której brzeg $S^{k-1}$
leży w już zbudowanej części.
Na
\[
 (D^k\times\{0\})\cup
 (S^{k-1}\times I)
\subset D^k\times I
\]
mamy dane: $f$ na dole
i homotopię $h$ na boku.
Ta podprzestrzeń jest
retraktem całego cylindra.
Można to zobaczyć jawnie:
pisząc $r=\|x\|$ i
\[
 \lambda(x,t)=
 \min\left\{\frac1r,\frac2{2-t}\right\},
 \qquad
 R(x,t)=\bigl(\lambda x,\,
                  2-\lambda(2-t)\bigr),
\]
przy $r=0$ przyjmujemy
$1/r=+\infty$. Dla $r\leq1/2$ minimum wynosi $2/(2-t)$,
więc wzór jest ciągły także przy $r=0$. Wartość
$R(x,t)$ leży albo na boku
$\|\lambda x\|=1$, albo
na dole z drugą współrzędną
równą $0$. Jeśli $t=0$
lub $r=1$, mamy $\lambda=1$,
więc $R$ pozostawia dane
punkty nieruchome.
Złożenie danych z $R$
przedłuża homotopię
na $D^k\times I$.

Dołączajmy teraz komórki
$X\setminus A$ według
wymiarów. Na każdym kroku
stosujemy poprzednią
konstrukcję na wszystkich
dyskach danego wymiaru;
na ich brzegach odwzorowania
już się zgadzają. Ponieważ
topologia CW jest zadana
przez odwzorowania charakterystyczne,
sklejone homotopie są
ciągłe. Dla nieskończenie wielu komórek uzasadnieniem ciągłości
z parametrem czasu jest Lemat~\ref{lem:zwartosc-cw-pi},
zastosowany do wszystkich map $D^k\times I\to Y$.
Dla $0$-komórek wybieramy homotopię stałą równą wartości $f$;
wzór na $R$ jest potrzebny dla $k\geq1$.
\end{proof}

\begin{twierdzenie}[Przybliżenie komórkowe]
\label{thm:przyblizenie-komorkowe}
Każde odwzorowanie $f:X\to Y$ kompleksów CW
jest homotopijne z odwzorowaniem
\emph{komórkowym}, czyli takim,
które spełnia
$f(X^k)\subseteq Y^k$ dla
każdego $k$. Jeśli na
podkompleksie $A\subseteq X$
odwzorowanie już jest komórkowe,
homotopię można wybrać
stałą na $A$.
\end{twierdzenie}
\begin{proof}
Wyjaśnijmy lokalny krok geometryczny. Niech
$u:D^k\to Y^m$ wysyła brzeg do $Y^{k-1}$, gdzie $m>k$.
Obraz jest zawarty w skończonym podkompleksie na mocy
Lematu~\ref{lem:zwartosc-cw-pi}. W jednej otwartej $m$-komórce
używamy współrzędnych $\operatorname{int}D^m$ i oznaczamy przez $v$
współrzędnościowy zapis $u$ na otwartym przeciwobrazie tej komórki.
Wybierzmy $0<a<b<c<1$. Zbiór $K=v^{-1}(\overline B_b)$
jest zwarty i leży we wnętrzu $D^k$: obraz $\overline B_b$
jest zwarty i domknięty w przestrzeni Hausdorffa $Y$.
Wybieramy skończone otoczenie wielościenne $P$ zbioru $K$
takie, że $K\subset\operatorname{int}P\subset P\subset v^{-1}(B_c)$.
Można je uzyskać jako sumę domkniętych kostek dostatecznie drobnej siatki
spotykających $K$. Triangulujemy $P$ i przez interpolację wartości $v$
w wierzchołkach dostajemy odcinkowo afiniczne $w$ z
$\|w-v\|<\varepsilon$, gdzie $\varepsilon<\min(b-a,1-c)$.
Istnienie takiego przybliżenia wynika z jednostajnej ciągłości $v$ na $P$:
na każdym małym sympleksie interpolujemy wartości różniące się o mniej niż $\varepsilon$.

Niech $\eta:D^k\to[0,1]$ będzie ciągła, równa $1$ na $K$
i $0$ poza $\operatorname{int}P$. Na $P$ stosujemy homotopię
$v_t=v+t\eta(w-v)$, a poza $P$ nic nie zmieniamy.
Wszystkie wartości pozostają we wnętrzu komórki i homotopia nie rusza brzegu dysku.
Poza $K$ mamy $\|v\|>b$, więc $\|v_t\|>a$.
W $K$ końcowe odwzorowanie jest równe $w$.
Jego obraz jest zawarty w skończonej sumie podprzestrzeni afinicznych
wymiaru co najwyżej $k<m$, która nie zawiera otwartej kuli $B_a$.
Istnieje więc $z\in B_a$ pominięty przez całe nowe odwzorowanie.

Promieniowa retrakcja $D^m\setminus\{z\}$ na $\partial D^m$
zostawia brzeg nieruchomy, więc po sklejeniu komórki daje homotopię
usuwającą jej wnętrze z obrazu $u$. Ciągłość z parametrem czasu
wynika z ilorazowego kryterium z poprzedniego lematu.
Usuwamy kolejno komórki najwyższego wymiaru i schodzimy z wymiarem do $k$.
Pracujemy w skończonym podkompleksie zawierającym obraz i domknięcia
spotkanych komórek, zatem proces kończy się po skończenie wielu krokach.

Zaczynamy od $X^0$:
obraz każdego wierzchołka
przesuwamy do odpowiedniego
wierzchołka w jego składowej
$Y$, a na $A$ niczego
nie zmieniamy. Załóżmy
indukcyjnie, że
$f(X^{k-1})\subseteq Y^{k-1}$.
Dla każdej $k$-komórki
stosujemy lokalny krok
do jej odwzorowania
charakterystycznego:
usuwamy po kolei obrazy
komórek celu o wymiarach
$m>k$, zachowując
jej brzeg. Dostajemy
odwzorowanie tej komórki do
$Y^k$. Dzięki
Lematowi~\ref{lem:hep-cw}
homotopię na dotychczas
zbudowanej części
przedłużamy na całe $X$
i przechodzimy do
następnego wymiaru.
Podane kroki wykonuje
się względnie $A$.
Etapowi wymiaru $k$ przydzielamy przedział
$[1-2^{-k},1-2^{-k-1}]$, a w chwili $1$ bierzemy odwzorowanie końcowe.
Każdy ustalony szkielet pozostaje nieruchomy po zakończeniu swojego etapu.
Na każdym $D^k\times I$ mamy więc skończenie wiele sklejonych homotopii
i stały odcinek końcowy. Lemat~\ref{lem:zwartosc-cw-pi}
zapewnia ciągłość całej homotopii, także w chwili $1$.
\end{proof}

\begin{uwaga}[Przybliżenie w równych wymiarach]
\label{rem:przyblizenie-rowne-wymiary}
Ta sama lokalna konstrukcja dla $k=m$ pozwala wybrać $z\in B_a$
poza obrazami ścian triangulacji $P$ i tych $m$-sympleksów,
na których $w$ ma rząd mniejszy od $m$. Jest to skończona suma
zbiorów o wymiarze mniejszym od $m$. Każdy przeciwobraz $z$
leży wtedy we wnętrzu jednego sympleksu, gdzie $w$ jest afinicznym homeomorfizmem.
Przeciwobrazów jest skończenie wiele, a każdy ma znak wyznacznika $+1$ lub $-1$.
Możemy wokół $z$ wybrać małą kulę, której przeciwobrazy są rozłącznymi
dyskami, odwzorowanymi homeomorficznie na tę kulę.
To wynika również ze zwartości: poza małymi otoczeniami znalezionych
przeciwobrazów obraz omija całe dostatecznie małe otoczenie $z$.
Ten wariant będzie użyty w dowodzie o pierwszych komórkach.
\end{uwaga}

\begin{wniosek}[Szkielet widziany przez $\pi_i$]
\label{cor:pi-szkielety}
Dla kompleksu CW $X$ para
$(X,X^n)$ jest $n$-spójna.
Włączenie $X^n\hookrightarrow X$
daje więc izomorfizm na
$\pi_i$ dla $i<n$
i surjekcję dla $i=n$.
Jeśli $1\leq i<k$, to
$\pi_i(S^k)=0$.
\end{wniosek}
\begin{proof}
Reprezentant klasy
$\pi_i(X,X^n)$ jest
odwzorowaniem $(D^i,\partial D^i)\to
(X,X^n)$, gdzie $i\leq n$.
Najpierw przybliżamy komórkowo mapę brzegu do $X^n$,
zachowując punkt bazowy, i przedłużamy tę homotopię na dysk
Lematem~\ref{lem:hep-cw}. Dopiero teraz brzeg jest komórkowy,
więc względne przybliżenie na dysku przesuwa cały obraz do
$X^i\subseteq X^n$. Klasa względna znika.
Każda składowa $X$ zawiera wierzchołek: punkt komórki można
połączyć z jej brzegiem i powtarzać to w malejących wymiarach.
To daje także stwierdzenie o składowych i przypadek $n=0$.
Długi ciąg~\eqref{eq:les-pi}
daje twierdzenie o
włączeniu szkieletu.
Sfera $S^k$ ma strukturę CW
z jednym wierzchołkiem
i jedną $k$-komórką,
więc jej $i$-szkielet
jest punktem dla $0\leq i<k$.
Każde oparte odwzorowanie
$S^i\to S^k$ jest
homotopijne do odwzorowania
komórkowego i wtedy
jego obraz jest punktem.
\end{proof}

\begin{stwierdzenie}[Wyższe grupy i nakrycia]
\label{prop:pi-n-nakrycie}
Jeśli $p:\widetilde X\to X$
jest nakryciem, $\widetilde x_0$
leży nad $x_0$, to
\[
 p_*:\pi_n(\widetilde X,\widetilde x_0)
 \xrightarrow{\ \cong\ }
 \pi_n(X,x_0)\qquad(n\geq2).
\]
\end{stwierdzenie}
\begin{proof}
Sfera $S^n$ dla $n\geq2$
jest jednospójna na mocy
Wniosku~\ref{cor:pi-szkielety}
(dla $n=2$ znika $\pi_1$,
a dla większych wymiarów
również). Każde oparte
odwzorowanie $S^n\to X$ podnosi
się więc jednoznacznie
po wskazaniu obrazu
punktu bazowego,
zgodnie z kryterium
z Twierdzenia~\ref{thm:kryterium-podnoszenia}.
Daje to surjektywność $p_*$.
Jeśli dwa podniesienia
po zrzutowaniu są
homotopijne przez odwzorowania oparte,
ich homotopię
$S^n\times I\to X$ podnosimy z kryterium podnoszenia:
$S^n\times I$ jest łukowo i lokalnie łukowo spójna oraz jednospójna,
bo retrahuje się na $S^n$. Ustalamy obraz $(\ast,0)$.
Ścieżka punktu bazowego leży w dyskretnym włóknie, więc jest stała,
a ograniczenie do $S^n\times\{0\}$ jest zadanym podniesieniem.
Jednoznaczność podnoszenia
sprawia, że na końcu
otrzymujemy drugie
podniesienie, ponieważ
oba zgadzają się w
punkcie bazowym.
To daje iniektywność.
\end{proof}

\begin{przyklad}[Okrąg i torus nie mają wyższych $\pi_n$]
\label{prz:pi-torus}
Nakrycia $\R\to S^1$ i
$\R^2\to T^2=S^1\times S^1$
mają ściągalne przestrzenie
całkowite. Zatem dla $n\geq2$
\[
 \pi_n(S^1)=0,\qquad
 \pi_n(T^2)=0.
\]
Jednocześnie
$H_2(T^2)=\Z$ z
Przykładu~\ref{prz:homologia-cw-torus}.
Już tu widać, że w
przestrzeni z nietrywialną
$\pi_1$ wyższa homologia
nie musi być obrazem
wyższej grupy homotopii.
\end{przyklad}

\section{Pierwsze odwzorowanie Hurewicza: od pętli do cykli}
\label{sec:hurewicz-jeden}

\begin{twierdzenie}[Hurewicz w stopniu pierwszym]
\label{thm:hurewicz-h1}
Dla przestrzeni łukowo spójnej $X$
odwzorowanie przypisujące opartej pętli
jej klasę singularnego
$1$-cyklu indukuje naturalny
izomorfizm
\[
 h_1:\pi_1(X,x_0)_{\mathrm{ab}}
 =\pi_1(X,x_0)/
 [\pi_1(X,x_0),\pi_1(X,x_0)]
 \xrightarrow{\ \cong\ }H_1(X;\Z).
\]
\end{twierdzenie}
\begin{proof}
Pętla $\ell:I\to X$ ma
$\partial\ell=[x_0]-[x_0]=0$.
Złożenie dwóch pętli
ma tę samą klasę homologii
co suma ich łańcuchów:
singularny $2$-sympleks, którego boki
przebiegają $\ell_1$, $\ell_2$ i ich złożenie, ma
brzeg $\ell_1+\ell_2-
(\ell_1*\ell_2)$
(z odpowiednią
parametryzacją boków).
Homotopia oparta pętli daje homologiczne cykle przez operator pryzmatyczny
z Twierdzenia~\ref{thm:homotopia-singularna}.
Zatem $[\ell]\mapsto
[\ell]_{H_1}$ jest
homomorfizmem do grupy
abelowej i zabija komutatory.

Skonstruujemy odwrotność.
Dla każdego $x\in X$
wybierzmy drogę $u_x$
od $x_0$ do $x$,
przy czym $u_{x_0}$
jest drogą stałą.
Singularnej krawędzi
$\sigma:x\to y$
przypiszmy w
$\pi_1(X,x_0)_{\mathrm{ab}}$
klasę pętli
\[
 L(\sigma)
 =[u_x*\sigma*\bar u_y].
\]
Niech $\tau:\Delta^2\to X$,
a $x_0',x_1',x_2'$
będą obrazami jej
wierzchołków.
Pętla złożona z boków
$\tau_{01}*\tau_{12}*
\bar\tau_{02}$ jest
zerowa w $\pi_1(X,x_0')$,
gdyż wypełnia ją
sam trójkąt $\tau$.
Po doprowadzeniu jej
do bazy drogami $u_{x_i'}$
otrzymujemy
\[
 L(\partial\tau)
 =L(\tau_{12})-L(\tau_{02})
  +L(\tau_{01})=0
 \quad\text{w }\pi_1(X,x_0)_{\mathrm{ab}}.
\]
Przez liniowość $L$
zeruje więc wszystkie
singularne brzegi
$2$-wymiarowe.
Na cyklach $z\in C_1(X)$
daje homomorfizm
$\bar L:H_1(X)\to
\pi_1(X,x_0)_{\mathrm{ab}}$.

Dla opartej pętli
$\ell$ mamy $L(\ell)=[\ell]$,
więc $\bar Lh_1=\id$.
W drugą stronę,
w homologii
obraz cyklu
$z=\sum m_\sigma\sigma$
rozwija się na
sumę łańcuchów $u_x$, samych krawędzi i $-u_y$,
z dokładnością do brzegów $2$-łańcuchów.
Pojedyncze $u_x$ zwykle nie są cyklami, więc nie przypisujemy
im osobnych klas w $H_1$. Rachunek przeprowadzamy w $C_1(X)/\operatorname{im}\partial_2$.
Tożsamość dla złożenia dróg wynika z trójkąta użytego na początku,
także gdy ich końce są różne; droga odwrotna daje ujemny łańcuch
modulo brzegi, bo droga złożona z własną odwrotnością jest homotopijna do stałej.
Współczynnik przy
każdej pomocniczej
drodze $u_x$ jest
równy minus współczynnikowi
przy $[x]$ w $\partial z$,
a zatem zerowy.
Pozostaje klasa
samego $z$ w $H_1(X)$;
formalnie używamy
homologii złożenia
dróg udowodnionej
na początku. Stąd
$h_1\bar L=\id$
i oba odwzorowania są
odwrotne.
\end{proof}

\begin{przyklad}[Od wolnej grupy do liczby pętli]
Bukiet $r$ okręgów ma
grupę podstawową
$F_r$ z
Przykładu~\ref{prz:pi1-dwa-okregi}
(dla $r=2$ rachunek
podano tam wprost).
Po narzuceniu
przemienności generatorów
$F_r^{\mathrm{ab}}\cong\Z^r$,
zgodnie z
Przykładem~\ref{prz:homologia-bukiet}.
Homologia pamięta tu
liczbę niezależnych
kierunków, ale nie
kolejność, w jakiej
przebywamy pętle.
\end{przyklad}

\section{Twierdzenie Whiteheada}
\label{sec:whitehead}

Równość wszystkich grup $\pi_n$
nie oznacza automatycznie, że
dowolne odwzorowanie przestrzeni
topologicznych ma homotopijną
odwrotność. Dla kompleksów CW
komórki pozwalają jednak
przekształcić informację
o sferach w deformację
całej przestrzeni.

\begin{definicja}[Słaba równoważność homotopijna]
\index{slaba rownowaznosc homotopijna@Słaba równoważność homotopijna}
\label{def:slaba-rownowaznosc}
Odwzorowanie $f:X\to Y$ jest
\emph{słabą równoważnością
homotopijną}, jeśli daje
bijekcję $\pi_0(X)\to\pi_0(Y)$
i dla wszystkich punktów
bazowych $x\in X$ oraz
wszystkich $n\geq1$
izomorfizmy
$f_*:\pi_n(X,x)\to
\pi_n(Y,f(x))$.
\end{definicja}

\begin{lemat}[Ściskanie komórek]
\label{lem:kompresja-cw}
Niech $(K,L)$ będzie parą CW,
a $B\subseteq Y$ niepustą
podprzestrzenią. Załóżmy, że
każda składowa łukowa $Y$ spotyka
$B$ oraz wszystkie grupy
$\pi_k(Y,B,b)$ dla
$k\geq1$ i $b\in B$
są trywialne (dla $k=1$ chodzi o zbiór z wyróżnioną klasą). Wówczas
każde odwzorowanie par
$F:(K,L)\to(Y,B)$
jest homotopijne
\emph{względem $L$}
z odwzorowaniem o obrazie w $B$.
\end{lemat}
\begin{proof}
Dla $0$-komórek $K\setminus L$
wybieramy drogi z ich obrazów
do $B$. Homotopia przesuwa
te obrazy do $B$ i jest
stała na $L$.

Załóżmy, że tak zrobiliśmy
na wszystkich komórkach
o wymiarze mniejszym od $k$.
Dla dołączonej $k$-komórki
odwzorowanie charakterystyczne
daje parę
$(D^k,S^{k-1})\to(Y,B)$.
Przy $k\geq1$ jej klasa
w $\pi_k(Y,B)$ jest zerowa.
Znaczy to, że można
deformować dysk do $B$,
pozwalając początkowo
brzegowi poruszać się w $B$.
Wyjaśnijmy, jak unieruchomić brzeg. Niech
$G:D^k\times I\to Y$ będzie względną homotopią do stałej,
z $G_0=F$ i $G_t(S^{k-1})\subset B$.
Połóżmy $a_t=1-t/2$ i dla $x=ru$ określmy
\[
 K_t(ru)=
 \begin{cases}
 G_t(ru/a_t),&r\leq a_t,\\
 G_{2(1-r)}(u),&a_t\leq r\leq1.
 \end{cases}
\]
Na $r=a_t$ oba wzory dają $G_t(u)$, na $r=1$ mamy stale $F(u)$,
a dla $t=0$ otrzymujemy $F$. Ciągłość przy zanikającym kołnierzu
wynika z ciągłości $G$; dla $t=1$ cały obraz leży w $B$.
Tak dostajemy deformację względem całego brzegu, a nie jedynie
homotopię, która pozwala mu przesuwać się w $B$.
Wykonujemy je na
wszystkich $k$-komórkach
równocześnie.
Lemat~\ref{lem:hep-cw}
przedłuża otrzymaną
homotopię do reszty
$K$, bez ruszania $L$.
Indukcja usuwa obrazy
wszystkich komórek
spoza $L$ z $Y\setminus B$.
Dla nieskończonego CW
przeznaczamy na etap $k$
coraz krótszy przedział
czasu kończący się w $1$.
Każdy ustalony szkielet staje się w końcu nieruchomy.
Ograniczenie do każdej domkniętej komórki, wraz z całym przedziałem czasu,
jest skończonym sklejeniem ciągłych homotopii i stałego końca.
Lemat~\ref{lem:zwartosc-cw-pi} daje ciągłość także w chwili $1$.
\end{proof}

\begin{twierdzenie}[Whitehead]
\label{thm:whitehead}
Słaba równoważność homotopijna
$f:X\to Y$ między kompleksami
CW jest równoważnością
homotopijną. Jeśli
$i:X\hookrightarrow Y$
jest włączeniem podkompleksu
i indukuje izomorfizmy
wszystkich $\pi_n$ oraz
bijekcję na $\pi_0$, to
$X$ jest silnym retraktem
deformacyjnym $Y$.
\end{twierdzenie}
\begin{proof}
Jeśli $X=\varnothing$, bijekcja na składowych wymusza $Y=\varnothing$
i teza jest natychmiastowa. Dalej zakładamy $X\ne\varnothing$.
Zacznijmy od włączenia.
W długim ciągu~\eqref{eq:les-pi}
pary $(Y,X)$ znajdują się
fragmenty
\[
 \pi_n(X)\xrightarrow[\cong]{i_*}
 \pi_n(Y)\longrightarrow
 \pi_n(Y,X)\longrightarrow
 \pi_{n-1}(X)
 \xrightarrow[\cong]{i_*}
 \pi_{n-1}(Y).
\]
Dokładność wymusza
$\pi_n(Y,X)=0$ dla
wszystkich $n\geq2$.
Dla stopnia $1$ warto użyć dróg, ponieważ nie jest to już ciąg grup.
Droga względna od $a\in X$ do $x_0$ ma oba końce w jednej składowej $X$,
gdyż $\pi_0(X)\to\pi_0(Y)$ jest injekcją. Dopinamy drogę
od $x_0$ do $a$ w $X$, otrzymując pętlę w $Y$.
Surjektywność $\pi_1(X)\to\pi_1(Y)$ pozwala zastąpić ją pętlą w $X$.
Każda taka droga jest trywialna względnie $X$, więc $\pi_1(Y,X)$
ma tylko wyróżniony element. Każda składowa $Y$ spotyka $X$,
bo odwzorowanie na $\pi_0$ jest także surjekcją.
Lemat~\ref{lem:kompresja-cw}
zastosowany do
$F=\id_Y:(Y,X)\to(Y,X)$
daje homotopię nieruchomą
na $X$, prowadzącą od
$\id_Y$ do odwzorowania
$r:Y\to X$.
Ponieważ $r|_X=\id_X$,
jest to silna retrakcja
deformacyjna.

Niech teraz $f$ będzie
dowolną słabą równoważnością
CW. Twierdzenie
\ref{thm:przyblizenie-komorkowe}
pozwala zastąpić ją
odwzorowaniem komórkowym $f'\simeq f$.
Zmiana punktu bazowego wzdłuż tej homotopii pokazuje,
że $f'$ nadal jest słabą równoważnością.
Utwórzmy jego
\emph{cylinder odwzorowania}
\[
 M_{f'}=(X\times I)\sqcup Y/
       \bigl((x,1)\sim f'(x)\bigr).
\]
Cylinder ma strukturę CW: poza komórkami $Y$ bierzemy dolne
komórki $e^k\times\{0\}$ i komórki $e^k\times(0,1)$ wymiaru $k+1$.
Brzeg każdej z tych ostatnich składa się z dołu, boków nad
brzegiem $e^k$ i góry odwzorowanej przez $f'$ do $Y^k$.
Komórkowość $f'$ zapewnia zgodność wymiarów, a topologia
ilorazowa i Lemat~\ref{lem:zwartosc-cw-pi} dają topologię słabą.
Dolna kopia $X=X\times\{0\}$ jest więc podkompleksem $M_{f'}$.
Zwijając każdy odcinek
$\{x\}\times I$ do
$f'(x)\in Y$, otrzymujemy
retrakcyjną deformację
$M_{f'}\to Y$.
Złożenie włączenia
$X\hookrightarrow M_{f'}$
z tą retrakcją to $f'$.
Ponieważ $f'_*$ i
retrakcyjne odwzorowanie na $Y$
indukują izomorfizmy
na $\pi_n$, również
włączenie $X\hookrightarrow
M_{f'}$ je indukuje.
Pierwsza część dowodu
ściska więc $M_{f'}$
do $X$. Oznaczmy dolne włączenie przez $j$, włączenie $Y$ do cylindra przez $i$,
a retrakcje na $X,Y$ odpowiednio przez $s,r$.
Wtedy $f'=rj$, a odwrotnością homotopijną jest $g=si$:
$gf'=sirj\simeq sj=\id_X$ oraz $f'g=rjsi\simeq ri=\id_Y$.
Ponieważ $f\simeq f'$, to samo $g$ jest odwrotnością homotopijną $f$.
\end{proof}

\begin{figure}[!htbp]
\centering
\begin{tikzpicture}[>=Stealth,line cap=round,line join=round]
 \fill[manifoldblue!10]
  (0,0)--(4.1,0)--(4.35,2.0)--(.25,2.0)--cycle;
 \draw[manifoldblue!70!black,thick]
  (0,0)--(4.1,0)--(4.35,2.0)--(.25,2.0)--cycle;
 \foreach \x in {.7,1.8,2.9,3.7}{
  \draw[manifoldblue!35] (\x,.07)--({\x+.23},1.94);
 }
 \fill[manifoldgreen!18]
  (-.38,2.0) rectangle (4.86,2.5);
 \draw[manifoldgreen!70!black,thick]
  (-.38,2.0) rectangle (4.86,2.5);
 \node at (2.2,2.25) {$Y$};
 \node[below] at (2.08,0) {$X\times\{0\}$};
 \node[manifoldblue!70!black] at (2.2,1.02)
  {$X\times I$};
 \node[manifoldred!85!black,above]
  at (2.2,2.62) {$(x,1)\sim f'(x)$};
 \draw[->,manifoldred,thick]
  (3.35,1.48)--(3.35,2.2)
  node[midway,right=3pt] {$r$};
 \node[font=\small,align=left,anchor=west,text width=4.8cm]
  at (5.15,1.15)
  {Cylinder zwija się na $Y$.
   Gdy $f'$ jest słabą równoważnością,
   poprzednia część dowodu zwija go
   również na dolną kopię $X$.};
\end{tikzpicture}
\caption{Cylinder odwzorowania zamienia
odwzorowanie $f':X\to Y$ na włączenie
podkompleksu $X\hookrightarrow M_{f'}$.
Rysunek pokazuje mechanizm dowodu
twierdzenia Whiteheada.}
\label{fig:cylinder-whiteheada}
\end{figure}

\begin{uwaga}[Zakres twierdzenia]
Sprawdzenie tylko $\pi_1$
i kilku pierwszych grup
$\pi_n$ nie wystarcza
do zastosowania Whiteheada:
założenie dotyczy
\emph{każdego} $n$.
Również samo podobieństwo
grup homologii dwóch
przestrzeni nie wskazuje
jeszcze odwzorowania, które
mogłoby być równoważnością.
Niżej wyprowadzimy
użyteczne kryterium
homologiczne przy
dodatkowej jednospójności.
\end{uwaga}

\section{Odwzorowanie Hurewicza i pierwszy niezerowy wymiar}
\label{sec:hurewicz-wyzszy}

\begin{definicja}[Odwzorowanie Hurewicza]
\index{odwzorowanie hurewicza@Odwzorowanie Hurewicza}
\label{def:mapa-hurewicza}
Homologia $H_n(S^n;\Z)$
jest izomorficzna z $\Z$
na mocy
Twierdzenia~\ref{thm:homologia-sfer}.
Wybierzmy generatory zgodnie z orientacjami dysków i ich brzegów,
tak aby $\delta[D^n,S^{n-1}]=[S^{n-1}]$.
Dla $n\geq2$ określamy
\[
 h_n:\pi_n(X,x_0)\longrightarrow H_n(X;\Z),
 \qquad h_n([f])=f_*[S^n].
\]
Dla pary $(X,A)$ bierzemy
analogicznie generator
$[D^n,S^{n-1}]
\in H_n(D^n,S^{n-1})$ i
kładziemy
\[
 h_n^{\mathrm{rel}}([F])
 =F_*[D^n,S^{n-1}]
 \in H_n(X,A).
\]
W stopniu $1$ dla przestrzeni używamy homomorfizmu
$\pi_1(X)\to\pi_1(X)_{\mathrm{ab}}\xrightarrow{h_1}H_1(X)$.
Względny homomorfizm określamy tutaj dla $n\geq2$;
zbiór $\pi_1(X,A)$ nie jest w ogólności grupą.
\end{definicja}

Homotopie oparte są w
szczególności zwykłymi
homotopiami odwzorowań, więc
$h_n$ jest dobrze
określone dzięki
Twierdzeniu~\ref{thm:homotopia-singularna}.
Odwzorowanie jest homomorfizmem:
sklejanie dwóch połówek
sfery odpowiada sumie
ich klas fundamentalnych.
Dokładniej, sklejanie jest złożeniem mapy ściskającej równik
$S^n\to S^n\vee S^n$ z $f\vee g$. Mapa ściskająca posyła
generator na sumę dwóch generatorów: rzut na każdy składnik
ma stopień $+1$, co widać z klasy orientacji odpowiedniej półkuli
i naturalności odwzorowania łączącego dla dysku i brzegu.
Względnie używamy kostki podzielonej na dwie połowy w pierwszej współrzędnej.
Po zorientowaniu obu połówek zgodnie z całą kostką ich wewnętrzne
ściany kasują się; obraz łańcucha fundamentalnego jest sumą obu obrazów.
To daje również homomorfizm względny. Naturalność wynika od razu ze składania map.

Do dowodu potrzebujemy
przeliczyć grupy homotopii
prostych modeli komórkowych.
Najważniejszy jest
następujący krok: w
pierwszym nowym wymiarze
komórki są niezależne,
a komórki o wymiar
wyższym wprowadzają
relacje.

\begin{lemat}[Pierwsze komórki i ich relacje]
\label{lem:pierwsze-komorki-pi}
Niech $A$ będzie jednospójnym
podkompleksem CW kompleksu $Z$.
Załóżmy, że wszystkie komórki
$Z\setminus A$ mają wymiar
co najmniej $n\geq2$.
Oznaczmy przez $I_n$
zbiór nowych $n$-komórek.
Wtedy
\[
 \pi_n(Z^n\cup A,A)
 \cong\bigoplus_{\alpha\in I_n}\Z[e_\alpha^n]
 \xrightarrow[\cong]{\ h_n^{\mathrm{rel}}\ }
 H_n(Z^n\cup A,A).
\]
Po dołączeniu $(n+1)$-komórek
obie grupy otrzymujemy
z tej wolnej grupy przez
te \emph{same} relacje:
obraz brzegu odwzorowania
dołączającego $S^n\to
Z^n\cup A$. Komórki
o wymiarach co najmniej
$n+2$ nie zmieniają
tych grup.
\end{lemat}
\begin{proof}
Niech $W=Z^n\cup A$ oraz $W'=Z^{n+1}\cup A$.
Przestrzeń $W$ jest jednospójna: do jednospójnego $A$
dołączamy tylko komórki wymiaru co najmniej $2$, co nie tworzy
nowych pętli (lokalny krok przybliżenia komórkowego przesuwa pętlę do $A$).
W szczególności także $\pi_2(W,A)$ jest abelowa, bo z ciągu pary
jest ilorazem abelowej $\pi_2(W)$ przy $\pi_1(A)=0$.
Możemy więc mówić o homomorfizmie z wolnej grupy abelowej
również w przypadku $n=2$.
Po zwinięciu $A$
otrzymujemy bukiet
$\bigvee_{\alpha\in I_n}S^n$.
Względna homologia
$H_n(W,A)$ jest
wolna na komórkach
na mocy
Twierdzenia~\ref{thm:warstwa-komorek}
i rozszerzenia na pary CW po Lemacie~\ref{lem:zwartosc-cw-pi};
to jeszcze nie jest
argument o homotopii,
więc rozpiszmy go
osobno.

Każdą mapę charakterystyczną nowej komórki opieramy w $x_0$
przez dołączenie drogi w $A$ do obrazu wybranego punktu jej brzegu.
Wybór drogi nie zmienia klasy, bo $A$ jest jednospójna.
Tak rozumiemy symbole $[e_\alpha^n]$ po stronie homotopii.
Weźmy teraz odwzorowanie względne
$F:(D^n,S^{n-1},*)\to(W,A,x_0)$.
Zwarty obraz spotyka
tylko skończenie wiele
otwartych komórek.
W każdej z nowych $n$-komórek stosujemy przybliżenie
z Uwagi~\ref{rem:przyblizenie-rowne-wymiary}, nieruchome przy $S^{n-1}$.
Wybieramy regularny punkt $z_\alpha$ poza obrazami ścian triangulacji.
Ma on skończenie wiele przeciwobrazów; na małych otoczeniach
$F$ jest afinicznym homeomorfizmem o znaku $+1$ lub $-1$.
Wybierzmy wokół $z_\alpha$ małą kulę. Jej przeciwobrazy są rozłącznymi
dyskami $D_j$, a poza nimi $F$ omija tę kulę.
Parametryzując $D_j$ przez odwrotność lokalnej mapy afinicznej,
otrzymujemy dokładnie lokalny fragment mapy charakterystycznej;
ujemny znak odpowiada odbiciu pierwszej współrzędnej, czyli odwrotności w grupie. Po usunięciu
środków każda komórka
promieniowo zwija się
na swój brzeg w $A$;
zatem odwzorowanie na dopełnieniu
małych dysków możemy
przesunąć do $A$.
Przedłużamy tę deformację na $D_j$ przez rozszerzanie ich obrazów
od małej kuli do całej komórki, zostawiając zgodność na wspólnych brzegach.
Połączmy małe dyski i punkt bazowy drzewem łuków w ich dopełnieniu.
Rozcinając wzdłuż cienkich pasów tego drzewa, przedstawiamy dysk
jako sumę dysków z doprowadzeniem do bazy; pozostałe części mają obraz w $A$
i reprezentują zero względne. Takie rozcięcie realizuje kolejno
kostkowe dodawanie z Definicji~\ref{def:pi-wzgledne}:
wewnętrzny pas ma obraz w $A$ i można go przesuwać podczas homotopii względnej.
Zmiana drogi doprowadzającej modyfikuje składnik przez pętlę w $A$.
Ponieważ $\pi_1(A)=0$, te zmiany nie wpływają na klasę.
Ostatecznie $[F]$ jest sumą klas charakterystycznych dysków
$e_\alpha^n$, z ich znakami. Każda suma jest skończona.
To dowodzi surjektywności
odwzorowania od wolnej grupy
na nowych komórkach
do $\pi_n(W,A)$.
Obrazy tych generatorów
przez $h_n^{\mathrm{rel}}$
tworzą bazę
$H_n(W,A)$. Gdyby ich
kombinacja dawała
zerową klasę względnej
homotopii, jej obraz
homologiczny również
byłby zerowy; niezależność
bazy wymusza wszystkie
współczynniki równe $0$.
Mamy więc izomorfizm.

Niech teraz dołączamy dyski $D^{n+1}_\beta$, otrzymując $W'$.
Każdą mapę $D^n\to W'$ względnie $A$ przesuwamy do $W$
przez lokalne usuwanie punktów w komórkach wymiaru $n+1$.
Włączenie indukuje więc surjekcję na $\pi_n$.
Relację od $\varphi_\beta:S^n\to W$ rozumiemy jako jej obraz
przez $\pi_n(W)\to\pi_n(W,A)$.
Ich odwzorowania charakterystyczne
dają oczywiste relacje:
brzeg każdego takiego
dysku staje się zerem
w $\pi_n(W',A)$.
Dlaczego nie pojawiają
się inne relacje?
Homotopię względną
między dwoma sumami
$n$-dysków przedstawiamy
na $D^n\times I$.
Uwaga~\ref{rem:przyblizenie-rowne-wymiary}, zastosowana do cylindra
triangulowanego jako $(n+1)$-dysk i nieruchomo na całym jego brzegu,
zapewnia, że przeciwobrazy wybranych punktów $(n+1)$-komórek
są skończonym zbiorem
we wnętrzu $(n+1)$-wymiarowego
cylindra. Usuwamy małe
otoczenia tych punktów.
Pozostała część odwzorowania
omija środki i promieniowo
przesuwa się do $W$.
Na brzegu każdego
usuniętego otoczenia
otrzymujemy, ze znakiem,
odwzorowanie dołączające
odpowiedniej komórki.
Łączymy te otoczenia z brzegiem cylindra drzewem łuków.
Po rozcięciu wzdłuż drzewa zorientowany brzeg pozostałego obszaru
daje różnicę końcowych $n$-dysków oraz sumę nowych sfer brzegowych.
Jego wypełnienie leży w $W$, a pozostałe ściany w $A$,
więc w $\pi_n(W,A)$ otrzymujemy równość różnicy klas z sumą
obrazów $[\varphi_\beta]$, ze znakami lokalnych homeomorfizmów.
Drogi doprowadzające leżą teraz w $W$; nie zmieniają klas,
ponieważ $W$ jest jednospójna. To wyjaśnia, dlaczego dostajemy
zwykłe całkowite kombinacje, bez dodatkowych sprzężeń.
Każda relacja w
$\pi_n(W',A)$
jest więc kombinacją
brzegów dołączonych
komórek.
W homologii to samo
stwierdza wzór
na brzeg komórkowy
\eqref{eq:formula-komorkowa}:
stopnie odwzorowań dołączających
są dokładnie
współczynnikami
tych relacji.

Wreszcie dla $k\geq n+2$
homotopię dysku $D^n$
można przybliżyć
komórkowo, zachowując
brzeg w $A$: jej obraz
nie musi wchodzić
do nowych $k$-komórek.
To samo dotyczy
$(n+1)$-wymiarowej
homotopii, bo $n+1<k$.
Zatem grupy $\pi_n$
nie zmieniają się.
Również względny
kompleks komórkowy
nie zmienia swoich
grup i brzegów
w stopniach $n$ i $n+1$.
\end{proof}

\begin{wniosek}[Zabijanie wybranych klas]
\label{cor:zabijanie-klas-pi}
Niech $Y$ będzie jednospójnym CW i $k\geq2$.
Dołączenie $(k+1)$-komórek wzdłuż klas $a_\beta\in\pi_k(Y)$
zmienia $\pi_k(Y)$ na iloraz przez podgrupę generowaną przez $a_\beta$,
a nie zmienia $\pi_i$ dla $i<k$.
\end{wniosek}
\begin{proof}
Oznaczmy nową przestrzeń przez $V$. Poprzedni lemat z bazą $A=Y$
i pierwszym nowym wymiarem $k+1$ mówi, że
$\pi_{k+1}(V,Y)$ jest wolna na dołączonych komórkach.
Mapa brzegowa do $\pi_k(Y)$ posyła ich generatory na $a_\beta$.
Niższe grupy względne znikają przez lokalny krok przybliżenia
komórkowego, ponieważ dyski tych wymiarów można przesunąć do $Y$.
Długi ciąg pary daje dokładnie podany iloraz i izomorfizmy w niższych stopniach.
\end{proof}

\begin{uwaga}[O przybliżeniu w lemacie]
Warunek $\pi_1(A)=0$ ma
konkretną rolę: pozwala
jednoznacznie przenosić
lokalne dyski do bazy.
Gdy $A$ nie jest
jednospójna, powstają
przesunięte przez pętle
kopie generatorów.
Wersja względnego
twierdzenia Hurewicza
bez tego warunku
wymaga ilorazu przez
działanie $\pi_1(A)$;
nie utożsamiamy wtedy
bezpośrednio całej
$\pi_n(X,A)$ z
$H_n(X,A)$.
\end{uwaga}

\begin{lemat}[Model bez niższych komórek]
\label{lem:model-wysoko-spojny}
Jeśli kompleks CW $X$
jest $(n-1)$-spójny,
$n\geq2$, to jest
homotopijnie równoważny
kompleksowi $Z$ z
jedną $0$-komórką
i bez komórek
wymiarów $1,\ldots,n-1$.
Odwzorowanie $Z\to X$
można dobrać tak,
by przenosiło
wyróżniony punkt
na $x_0$.
\end{lemat}
\begin{proof}
Zaczynamy od jednej
$0$-komórki $Z^0=\{*\}$
wysłanej do $x_0$.
Ponieważ $\pi_i(X)=0$
dla $i<n$, odwzorowanie to
jest już izomorfizmem
na tych grupach.
Dołączamy do punktu
po jednej $n$-komórce
dla wybranego zbioru
generatorów $\pi_n(X)$.
Każda taka komórka
z brzegiem sklejonym
do punktu jest $S^n$;
odwzorowanie do $X$ definiujemy
na niej reprezentantem
wybranej klasy.
Uzyskujemy surjekcję
na $\pi_n(X)$.

Teraz dla każdego
generatora jądra
$\pi_n(Z^n)\to\pi_n(X)$
dołączamy $(n+1)$-komórkę
wzdłuż reprezentującego
go odwzorowania $S^n\to Z^n$.
Jej obraz w $X$
jest zerowy w $\pi_n(X)$,
więc wybieramy
wypełnienie $D^{n+1}\to X$
i przedłużamy odwzorowanie
$Z\to X$.
Wniosek~\ref{cor:zabijanie-klas-pi}, zastosowany do $Y=Z^n$,
mówi, że dokładnie te relacje zabijają wybrane jądro; w
stopniu $n$ odwzorowanie
staje się izomorfizmem.
Dołączone komórki
nie zmieniają niższych
grup dzięki
Wnioskowi~\ref{cor:pi-szkielety}.

Powtarzamy procedurę
w stopniu $k=n+1,n+2,\ldots$:
komórki $k$-wymiarowe
dołączone do punktu
realizują brakujące
elementy $\pi_k(X)$,
a komórki $(k+1)$-wymiarowe
dołączone wzdłuż klas
z jądra zabijają
dokładnie te klasy.
Dołączenie sfer $S^k$ do punktu daje surjekcję na brakujące klasy
i nie narusza już ustalonych niższych grup: nowy bukiet ma retrakcję
na poprzedni etap, a względne dyski wymiaru $<k$ można przesunąć do niego.
Do zabijania jądra w każdym stopniu stosujemy
Wniosek~\ref{cor:zabijanie-klas-pi} z bazą równą \emph{całemu bieżącemu etapowi}.
Nie wolno tu ponownie traktować punktu jako jedynej bazy pierwszych komórek,
bo wcześniej dołączono już komórki niższych wymiarów.
Mapy dołączające przybliżamy komórkowo, dzięki czemu otrzymujemy strukturę CW;
homotopia takiej mapy nie zmienia jej klasy ani możliwości wypełnienia w $X$.
Wszystkie odwzorowania do $X$ przedłużają się,
bo obrazy odwzorowań
dołączających z jądra
są zerowe w $\pi_k(X)$.
Komórki późniejszego
wymiaru nie naruszają
ustalonych już
izomorfizmów w
niższych stopniach.
Po wykonaniu
wszystkich etapów otrzymujemy CW $Z$.
Mapy sfer i homotopie mają zwarte dziedziny, więc z Lematu~\ref{lem:zwartosc-cw-pi}
każda klasa i każda relacja pojawiają się już w skończonym etapie.
Stabilizacja grup oznacza zatem, że odwzorowanie $Z\to X$
indukuje izomorfizmy
na wszystkich
$\pi_k$ oraz na
$\pi_0$. Twierdzenie
\ref{thm:whitehead}
zamienia tę słabą
równoważność CW
w równoważność
homotopijną.
\end{proof}

\begin{twierdzenie}[Hurewicz: pierwszy niezerowy wymiar]
\label{thm:hurewicz-wyzszy}
Jeżeli $X$ jest
$(n-1)$-spójnym
kompleksem CW,
$n\geq2$, to
\[
 \widetilde H_i(X;\Z)=0
 \quad\text{dla }i<n,
 \qquad
 h_n:\pi_n(X,x_0)
 \xrightarrow{\ \cong\ }H_n(X;\Z).
\]
Izomorfizm jest naturalny:
dla odwzorowania opartego
$f:X\to Y$ odpowiedni
kwadrat z $f_*$ na
homotopii i homologii
jest przemienny.
\end{twierdzenie}
\begin{proof}
Lemat~\ref{lem:model-wysoko-spojny}
pozwala zastąpić
$X$ homotopijnie
równoważnym $Z$
z $Z^{n-1}=\{*\}$.
Homologia i homotopia
nie zmieniają się
przy tej operacji
(Twierdzenia
\ref{thm:homotopia-singularna}
i~\ref{thm:whitehead}).
W kompleksie komórkowym
$Z$ nie ma grup
łańcuchów w
stopniach $1,\ldots,n-1$,
więc
$\widetilde H_i(Z)=0$
dla $i<n$.

Niech
$E_n=\bigoplus_{\alpha\in I_n}
\Z[e_\alpha^n]$.
Z Lematu~\ref{lem:pierwsze-komorki-pi}
dla $A=\{*\}$ mamy
$\pi_n(Z^n)\cong E_n$.
Włączenie
$Z^n\hookrightarrow Z$
jest surjekcją
na $\pi_n$,
a jądrem są relacje
narzucone przez
$(n+1)$-komórki.
Po lewej otrzymujemy
zatem
\[
 \pi_n(Z)
 \cong
 E_n/\langle
   [\varphi_\beta]:
   \varphi_\beta:S^n\to Z^n
   \text{ dołącza }e_\beta^{n+1}
 \rangle.
\]
Po prawej
$\partial_n^{\mathrm{CW}}=0$,
bo $C_{n-1}^{\mathrm{CW}}=0$,
więc
\[
 H_n(Z)
 \cong E_n/
     \operatorname{im}
       \partial_{n+1}^{\mathrm{CW}}.
\]
Obraz $[\varphi_\beta]$
przez $h_n$ jest
klasą brzegową
$(n+1)$-komórki:
jej współczynniki
w bazie $E_n$
to te same stopnie
odwzorowań dołączających
z Twierdzenia
\ref{thm:formula-komorkowa}.
Oba ilorazy są
więc ilorazami
\emph{tej samej}
wolnej grupy przez
\emph{ten sam}
podmoduł.
Ponieważ $h_n$
na generatorach
$[e_\alpha^n]$
przypisuje im
klasy komórkowe
$[e_\alpha^n]$,
jest izomorfizmem.

Na koniec
$h_n(f_*[u])
=(f\circ u)_*[S^n]
=f_*(u_*[S^n])$,
co dowodzi
naturalności.
\end{proof}

\begin{twierdzenie}[Hurewicz względny dla par jednospójnych]
\label{thm:hurewicz-wzgledny}
Niech $(X,A)$ będzie parą CW,
$A\ne\varnothing$ i $A$
jednospójna. Jeśli para
jest $(n-1)$-spójna,
$n\geq2$, to
\[
 H_i(X,A;\Z)=0\quad(i<n),
 \qquad
 h_n^{\mathrm{rel}}:
 \pi_n(X,A,x_0)
 \xrightarrow{\ \cong\ }
 H_n(X,A;\Z).
\]
\end{twierdzenie}
\begin{proof}
Zbudujemy najpierw
model pary bez nowych
komórek poniżej
wymiaru $n$.
Długi ciąg~\eqref{eq:les-pi}
oraz $\pi_i(X,A)=0$
dają izomorfizmy
$\pi_i(A)\to\pi_i(X)$
dla $i<n-1$ i
surjekcję dla
$i=n-1$.
Zaczynamy od
identyczności $A\to A$
oraz włączenia $A\to X$.
Dołączamy $n$-komórki
wzdłuż reprezentantów
jądra
$\pi_{n-1}(A)\to
\pi_{n-1}(X)$.
Każdy taki reprezentant
wypełnia się w $X$,
więc odwzorowanie do $X$
przedłuża się na
nową komórkę.
W długim ciągu
pary nowych komórek
odwzorowanie brzegowe
$\pi_n(Z^n,A)\to
\pi_{n-1}(A)$
posyła każdą
charakterystyczną
$n$-komórkę na
klasę jej odwzorowania
dołączającego.
Lemat~\ref{lem:pierwsze-komorki-pi},
zastosowany do
jednospójnego $A$,
pokazuje, że
te dołączenia
zabijają dokładnie
podgrupę przez
nie generowaną,
czyli wybrane jądro.
W stopniach
niższych niczego
nie zmieniają.
Gdy $n=2$,
wyjściowe $\pi_1(A)$
jest już zerowe
i ten krok nie
wymaga komórek.

Teraz dla $k=n,n+1,\ldots$
postępujemy jak w
dowodzie Lematu
\ref{lem:model-wysoko-spojny}:
dołączamy $k$-komórki
z brzegiem wysłanym
do $x_0$, aby
otrzymać surjekcję
na $\pi_k(X)$,
a potem
$(k+1)$-komórki
wzdłuż generatorów
jądra, aby
otrzymać izomorfizm.
Nigdy nie dołączamy
komórek wymiaru
mniejszego od $n$.
W każdym kroku wniosek o zabijaniu klas stosujemy z bieżącym,
jednospójnym etapem jako bazą. Przybliżamy komórkowo mapy dołączające.
Powstaje odwzorowanie par $F:(Z,A)\to(X,A)$
równe identyczności na $A$, które
indukuje izomorfizmy
wszystkich grup
homotopii przestrzeni
oraz składowych.
Twierdzenie~\ref{thm:whitehead}
zapewnia równoważność
homotopijną $Z\simeq X$.
Porównując długie ciągi
homotopii i homologii
obu par, widzimy
ponadto, że $F$
indukuje izomorfizmy
grup \emph{względnych}:
odwzorowania na wyrazach dotyczących $A$ są identycznością,
a na wyrazach dotyczących $Z,X$ izomorfizmami.
Dla stopni co najmniej $3$ stosujemy lemat pięciu do grup abelowych.
W stopniu $2$ względne grupy są ilorazami $\pi_2(Z)$ i $\pi_2(X)$
przez obrazy $\pi_2(A)$, ponieważ $\pi_1(A)=0$;
to daje izomorfizm bez stosowania lematu pięciu do zbiorów.
Stopień $1$ jest trywialny po obu stronach, bo wszystkie trzy przestrzenie
są jednospójne. Dla homologii stosujemy zwykły lemat pięciu.

W modelu $(Z,A)$
nie ma względnych
łańcuchów komórkowych
w wymiarach $<n$,
więc $H_i(Z,A)=0$
dla $i<n$.
Lemat~\ref{lem:pierwsze-komorki-pi}
oblicza zarówno
$\pi_n(Z,A)$,
jak i $H_n(Z,A)$
jako tę samą
wolną grupę na
$n$-komórkach
modulo te same
relacje od
$(n+1)$-komórek.
Na generatorach
odwzorowanie $h_n^{\mathrm{rel}}$
przypisuje dyskowi
jego klasę orientacji,
więc jest izomorfizmem.
Naturalność
$h_n^{\mathrm{rel}}$
przenosi ten wynik
przez $F$
na parę $(X,A)$.
\end{proof}

\begin{przyklad}[Dysk względem brzegu]
\label{prz:pi-dysk-brzeg}
Dla pary $(D^{n+1},S^n)$
długi ciąg homotopii
i ściągalność dysku
dają w szczególności
\[
 \partial:
 \pi_{n+1}(D^{n+1},S^n)
 \xrightarrow{\ \cong\ }
 \pi_n(S^n)
 \qquad(n\geq1).
\]
Generator po lewej jest reprezentowany przez dysk tożsamościowy,
a jego brzegiem jest
tożsamość sfery.
W homologii dzieje się
to samo na mocy
Twierdzenia~\ref{thm:homologia-sfer};
oba ciągi łączy odwzorowanie
Hurewicza.
\end{przyklad}

\begin{wniosek}[Homologiczne kryterium Whiteheada]
\label{cor:whitehead-homologiczny}
Niech $f:X\to Y$ będzie odwzorowaniem
między jednospójnymi
kompleksami CW.
Jeśli
$f_*:H_i(X;\Z)\to H_i(Y;\Z)$
jest izomorfizmem
dla wszystkich $i$,
to $f$ jest
równoważnością
homotopijną.
\end{wniosek}
\begin{proof}
Zastępujemy $f$
odwzorowaniem komórkowym
homotopijnym z nim
i tworzymy cylinder
$M_f$ z dowodu
Twierdzenia~\ref{thm:whitehead}.
Włączenie
$i:X\hookrightarrow M_f$
ma te same odwzorowania
na homologii co $f$
po złożeniu z
retrakcyjną deformacją
$M_f\to Y$.
Zatem w długim
ciągu homologii
pary $(M_f,X)$
wszystkie grupy
$H_i(M_f,X)$
są zerowe.
Ponieważ $X$
i $M_f\simeq Y$
są jednospójne,
długi ciąg homotopii
daje
$\pi_1(M_f,X)=0$.

Załóżmy, że pewna
wyższa względna
grupa homotopii
nie znika, i weźmy
najmniejsze $n\geq2$
z $\pi_n(M_f,X)\ne0$.
Para jest wtedy
$(n-1)$-spójna,
a $X$ jednospójna.
Twierdzenie
\ref{thm:hurewicz-wzgledny}
dałoby
$H_n(M_f,X)\cong
\pi_n(M_f,X)\ne0$,
sprzeczność.
Wszystkie względne
$\pi_n$ są więc
zerowe, a długi
ciąg homotopii
pokazuje, że
$i$, a zatem $f$,
indukuje izomorfizmy
wszystkich $\pi_n$.
Twierdzenie
\ref{thm:whitehead}
kończy dowód.
\end{proof}

\section{Obliczenia i ograniczenia twierdzeń}
\label{sec:pi-przyklady}

\begin{przyklad}[Pierwsza grupa homotopii sfery]
\label{prz:pi-sfera-hurewicz}
Z Wniosku~\ref{cor:pi-szkielety}
wiemy, że $S^n$
jest $(n-1)$-spójna.
Twierdzenia
\ref{thm:hurewicz-wyzszy}
i~\ref{thm:homologia-sfer}
dają dla $n\geq2$
\[
 \pi_i(S^n)=0\quad(1\leq i<n),
 \qquad
 \pi_n(S^n)\cong\Z.
\]
Tożsamość $S^n\to S^n$
odpowiada generatorowi,
a odwzorowanie $g:S^n\to S^n$
odpowiada liczbie
$\deg g$ z
Definicji~\ref{def:stopien-sfer}.
Zatem oparte $g$ jest homotopijne do tożsamości lub do ustalonego
odbicia sfery zachowującego punkt bazowy odpowiednio wtedy,
gdy stopień wynosi $+1$ lub $-1$.
Odbicie ma stopień $-1$, bo odwraca orientację dysku i jego brzegu.
Dla nieopartego $g$ najpierw przesuwamy obraz wyróżnionego punktu do bazy
i przedłużamy tę homotopię na sferę przez własność HEP.
Jednospójność sfery sprawia, że wybór drogi nie wpływa na wynik. W szczególności
odwzorowanie sfery stopnia
$\pm1$ jest
równoważnością
homotopijną:
wywołuje izomorfizm
na $H_n$ i $H_0$,
a inne grupy homologii
sfery są zerowe,
więc można użyć
Wniosku~\ref{cor:whitehead-homologiczny}.
\end{przyklad}

\begin{przyklad}[Kompleks Moore'a]
\label{prz:moore-hurewicz}
Niech $n\geq2$ i $m\geq2$.
Do $S^n$ dołączmy
jedną $(n+1)$-komórkę
odwzorowaniem stopnia $m$.
Powstała przestrzeń
$M(\Z/m\Z,n)$ ma
kompleks komórkowy
\[
 0\longrightarrow\Z
 \xrightarrow{\;\times m\;}\Z
 \longrightarrow0
\]
w dodatnich stopniach $n+1$ i $n$; ponadto $C_0=\Z$,
a pozostałe grupy łańcuchów są zerowe. Zatem
$H_n(M)=\Z/m\Z$,
$H_{n+1}(M)=0$.
Ponieważ przestrzeń
jest $(n-1)$-spójna,
Hurewicz daje
$\pi_n(M)=\Z/m\Z$.
Intuicja jest
dosłowna: dołączony
dysk czyni
$m$-krotne
owinięcie sfery
zerem w pierwszej
niezerowej grupie.
\end{przyklad}

\begin{przyklad}[Dlaczego potrzeba jednospójności]
\label{prz:rp2-pi2}
Nakrycie antypodalne
$p:S^2\to\mathbb{RP}^2$
indukuje izomorfizm
$\pi_2(S^2)\cong
\pi_2(\mathbb{RP}^2)$
na mocy
Stwierdzenia~\ref{prop:pi-n-nakrycie}.
Zatem
$\pi_2(\mathbb{RP}^2)\cong\Z$.
Tymczasem
$H_2(\mathbb{RP}^2)=0$
z Przykładu~\ref{prz:homologia-rp2},
więc tutaj
$h_2:\Z\to0$
nie jest izomorfizmem.
Przestrzeń rzutowa
nie jest jednospójna:
$\pi_1(\mathbb{RP}^2)=\Z/2\Z$.
To konkretna granica
Twierdzenia~\ref{thm:hurewicz-wyzszy}.
\end{przyklad}

\begin{przyklad}[Bukiet $S^1\vee S^2$]
\label{prz:pi-bukiet-pokrycie}
Przestrzeń
$X=S^1\vee S^2$
ma $\pi_1(X)\cong\Z$.
Jej nakrycie uniwersalne
składa się z prostej
z kopią $S^2$ w każdym
punkcie całkowitym.
Pokrycie otrzymujemy przez przesunięcia prostej o liczby całkowite,
przenoszące jednocześnie $S_k^2$ na $S_{k+1}^2$.
Jego przestrzeń jest jednospójna: każda pętla ma obraz
w skończonym odcinku z dołączonymi skończenie wieloma sferami,
a taka przestrzeń jest jednospójna z van Kampena.
Można obliczyć homologię bez zwijania całej prostej:
mamy wierzchołki w punktach całkowitych, krawędzie między nimi
i po jednej $2$-komórce każdej sfery, dołączonej stale do jej wierzchołka.
Z homologii komórkowej
\[
 H_2(\widetilde X)
 \cong\bigoplus_{k\in\Z}\Z.
\]
Nakrycie nie zmienia
$\pi_2$, a
$\widetilde X$
jest jednospójna,
więc Hurewicz daje
\[
 \pi_2(X)\cong
 \pi_2(\widetilde X)
 \cong\bigoplus_{k\in\Z}\Z,
 \qquad
 H_2(X)\cong\Z.
\]
Każda podniesiona
sfera $S_k^2$
rzutuje się
na tę samą sferę
wewnątrz $X$.
W tych bazach
odwzorowanie Hurewicza
sumuje współczynniki:
$(a_k)_{k\in\Z}
\mapsto\sum_k a_k$.
Homologia traci
informację, którą
z podniesionych
sfer reprezentuje
dana klasa.
\end{przyklad}

\begin{figure}[!htbp]
\centering
\begin{tikzpicture}[>=Stealth,line cap=round,line join=round]
 \draw[manifoldgold!80!black,line width=2pt]
  (-4.5,0)--(4.5,0);
 \foreach \x/\lab in {-3/{-1},0/{0},3/{1}}{
  \shade[inner color=white,outer color=manifoldblue!37]
   (\x,.95) circle(.95);
  \draw[manifoldblue!75!black,thick]
   (\x,.95) circle(.95);
  \draw[manifoldblue!45,densely dashed,thin]
   ({\x+.95},.95)
   arc[start angle=0,end angle=180,x radius=.95,y radius=.23];
  \draw[manifoldblue!70!black,thin]
   ({\x-.95},.95)
   arc[start angle=180,end angle=360,x radius=.95,y radius=.23];
  \draw[manifoldblue!32,densely dashed,thin]
   (\x,1.9)..controls({\x-.28},1.6)
     and({\x-.28},.5)..(\x,0);
  \draw[manifoldblue!62!black,thin]
   (\x,1.9)..controls({\x+.28},1.6)
     and({\x+.28},.5)..(\x,0);
  \fill[manifoldred] (\x,0) circle(2.8pt);
  \node[below=5pt,font=\small] at (\x,0)
   {$k=\lab$};
 }
 \node[manifoldgold!75!black,font=\small]
  at (4.12,-.35) {$\R$};
 \node[font=\large] at (-4.28,1.02) {$\cdots$};
 \node[font=\large] at (4.28,1.02) {$\cdots$};
\end{tikzpicture}
\caption{Fragment nakrycia uniwersalnego
$S^1\vee S^2$: kopia sfery w każdym
punkcie całkowitym nad pętlą okręgu.
Każda sfera daje inny generator
$\pi_2(X)$, ale wszystkie mają
ten sam obraz w $H_2(X)$.}
\label{fig:nakrycie-bukiet-sfer}
\end{figure}

\begin{uwaga}[Który Whitehead?]
Twierdzenie~\ref{thm:whitehead}
dotyczy słabych równoważności
homotopijnych CW.
\emph{Torsja Whiteheada}
jest innym, późniejszym
niezmiennikiem używanym
przy prostej homotopii
i ogólniejszych
twierdzeniach o
kobordyzmach; sama
nazwa nie czyni jej
częścią tego twierdzenia.
\end{uwaga}

\section*{Dalsza droga}
\addcontentsline{toc}{section}{Dalsza droga}
Następny naturalny blok to
formy różniczkowe,
orientacja, całkowanie,
twierdzenie Stokesa
i kohomologia de Rhama.
Po nich rozwiniemy
geometrię riemannowską:
tensor metryczny, koneksję,
symbole Christoffela,
transport równoległy,
geodezyjne, odwzorowanie
wykładnicze, twierdzenie
Hopfa--Rinowa, krzywiznę
i tensory krzywizny,
pola Jacobiego,
twierdzenia porównawcze
oraz twierdzenie
Gaussa--Bonneta.
\emph{Dopiero po tym bloku}
przejdziemy do teorii
Morse'a, uchwytów,
twierdzenia Smale'a
i chirurgii.
Nie wprowadzamy
tych późniejszych
teorii w obecnym rozdziale.

\section*{Dalsza lektura}
\addcontentsline{toc}{section}{Dalsza lektura}
Allen Hatcher,
\emph{Algebraic Topology},
rozdział 4, §§4.1--4.2,
\url{https://pi.math.cornell.edu/~hatcher/AT/ATch4.pdf},
zawiera pełne rozwinięcie
grup względnych,
przybliżenia komórkowego,
twierdzeń Whiteheada
i Hurewicza.
\clearpage
